\section{Convergence and interchange of expectation}\label{sec:convergence}

\subsection{Why pointwise convergence is insufficient}

\begin{example}[Moving spike]\label{ex:moving-spike}
Let the countable sample space be \(\N\) with
\(\Pp(\{k\})=2^{-k-1}\).  Define
\[
 X_n(k)=2^{n+1}\ind_{k=n}.
\]
For every fixed \(k\), \(X_n(k)=0\) for all \(n>k\), so \(X_n\to0\)
pointwise.  Nevertheless,
\[
 \E[X_n]=2^{-n-1}2^{n+1}=1
\]
for every \(n\).  Hence pointwise convergence alone does not imply convergence
of expectations.
\end{example}

This counterexample is fully effective: the geometric law has an exact tail,
the values are natural numbers, and the pointwise convergence modulus at
outcome \(k\) is \(k+1\).  What fails is a uniform integrability bound.

\begin{example}[Interchanging an infinite sum and expectation can fail]
\label{ex:series-interchange}
Put \(X_{-1}=0\), \(Y_0=X_0\), and
\(Y_n=X_n-X_{n-1}\) for \(n\ge1\), using the moving spikes above.  Then
\[
 \sum_{j=0}^{N}Y_j=X_N\longrightarrow0
\]
pointwise, so the pointwise series sum is zero.  But
\[
 \sum_{j=0}^{N}\E[Y_j]=\E[X_N]=1
\]
for every \(N\).  Thus
\(\E[\sum_jY_j]=0\ne1=\sum_j\E[Y_j]\).
Absolute summability or domination hypotheses are essential.
\end{example}

\subsection{Effective monotone convergence}

\begin{theorem}[Monotone convergence as a lower-real identity]
\label{thm:monotone-convergence}
Let \(0\le X_0\le X_1\le\cdots\) be a primitive-recursive sequence of
nonnegative finite simple-law names on compatible refinements.  Define the
pointwise supremum by its increasing approximants.  Then
\[
 \E[\sup_nX_n]=\sup_n\E[X_n]
\]
as equality of lower reals.  If an expectation-tail modulus is supplied, both
sides are equal as regular reals.
\end{theorem}

\begin{proof}
At every finite stage, expectation preserves order, so
\(\E[X_n]\) is increasing and bounded above by every finite approximation to
\(\E[\sup X_n]\).  Conversely, the definition of the supremum name uses
\(X_n\) itself as its increasing simple approximants; therefore its lower-real
expectation is, by construction, the directed family of the finite
expectations.  More explicitly, any rational lower bound certified for one
side occurs at a finite level and is certified for the other side at the same
level.  If a modulus says the omitted expectation after stage \(N(p)\) is at
most \(2^{-p}\), the common lower real becomes a regular interval name.
\end{proof}

\begin{corollary}[Nonnegative series and Tonelli form]
\label{cor:tonelli}
For a primitive-recursive sequence of nonnegative observables,
\[
 \E\left[\sum_{n=0}^{\infty}X_n\right]
 =\sum_{n=0}^{\infty}\E[X_n]
\]
as lower reals.  With an absolute expectation-tail modulus, the equality is
one of regular reals and the sum may be effectively regrouped.
\end{corollary}

\begin{proof}
Apply \Cref{thm:monotone-convergence} to the finite partial sums and use exact
finite linearity at every stage.
\end{proof}

\subsection{Finite Fatou kernel and effective liminf}

\begin{lemma}[Finite Fatou kernel]\label{lem:finite-fatou}
For nonnegative finite observables \(X_N,\ldots,X_M\),
\[
 \E\left[\min_{N\le n\le M}X_n\right]
 \le \min_{N\le n\le M}\E[X_n].
\]
\end{lemma}

\begin{proof}
The pointwise minimum is at most each \(X_n\).  Monotonicity gives an
expectation inequality for each \(n\); bounded search takes the least rational
right side.  The minimum evaluator and the selected index are primitive
recursive on the finite list.
\end{proof}

\begin{theorem}[Effective Fatou lemma]\label{thm:fatou}
Let \(X_n\ge0\) be an effective sequence.  For every \(N\), suppose the
finite minima
\[
 Y_{N,M}=\min_{N\le n\le M}X_n
\]
are supplied on compatible simple-law refinements and converge, as
\(M\to\infty\), to a name \(Y_N\) with an explicit \(L^1\) modulus.  Suppose
also that the increasing sequence \(Y_N\) presents the pointwise
\(\liminf\).  Then
\[
 \E[\liminfN_{n\to\infty}X_n]
 \le\liminfN_{n\to\infty}\E[X_n]
\]
as a lower-real inequality; with regular moduli it is an inequality of regular
reals.
\end{theorem}

\begin{proof}
For finite \(M\), \Cref{lem:finite-fatou} gives
\[
 \E[Y_{N,M}]
 \le\min_{N\le n\le M}\E[X_n].
\]
The supplied \(L^1\) modulus and
\(\lvert\E U-\E V\rvert\le\E\lvert U-V\rvert\) let \(M\) increase on the
left, yielding \(\E[Y_N]\).  The finite minima on the right decrease to the
upper-real tail infimum \(\inf_{n\ge N}\E[X_n]\); a regular version uses the
corresponding numerical modulus.  Now \(Y_N\) increases with \(N\), so
effective monotone convergence gives the expectation of the liminf.  Taking
the increasing supremum over \(N\) on the right gives the numerical liminf.
At each requested precision only finite minima, one finite expectation theorem,
and the stated moduli are used.
\end{proof}

\subsection{Uniform integrability and dominated convergence}

\begin{definition}[Effective uniform integrability]\label{def:ui}
A sequence \((X_n)\) is effectively uniformly integrable when a
primitive-recursive threshold \(K(p)\) satisfies
\[
 \E\bigl[\abs{X_n}\ind_{\abs{X_n}>K(p)}\bigr]\le2^{-p}
 \quad\text{for every }n.
\]
The universal quantifier is implemented by a proof schema for the fixed
sequence evaluator, not by bounded search over all \(n\).
\end{definition}

\begin{lemma}[Truncation estimate]\label{lem:truncation}
Suppose \(\abs X,\abs Y\le G\), where \(G\ge0\).  For rational
\(K,\eta>0\),
\[
 \E\abs{X-Y}
 \le \eta+2K\Pp(\abs{X-Y}>\eta)
      +2\E[G\ind_{G>K}].
\]
\end{lemma}

\begin{proof}
Split the sample space into three decidable regions.  On
\(\abs{X-Y}\le\eta\), the contribution is at most \(\eta\).  On the remaining
region with \(G\le K\), we have \(\abs{X-Y}\le\abs X+\abs Y\le2K\), giving the
middle term.  On \(G>K\), use \(\abs{X-Y}\le2G\).  Indicator linearity and the
finite union bound produce the displayed inequality.  For effective names,
apply it on a sufficiently fine common finite refinement and absorb the
\(L^1\) approximation errors.
\end{proof}

\begin{theorem}[Effective dominated convergence]\label{thm:dct}
Let \(X_n,X,G\) be effective \(L^1\) names satisfying
\(\abs{X_n},\abs X\le G\).  Assume:
\begin{enumerate}
\item an integrability-tail modulus for \(G\), and
\item a convergence-in-probability modulus: given \(p,q\), a
  primitive-recursive \(N(p,q)\) makes
  \[
   n\ge N(p,q)\Longrightarrow
   \Pp(\abs{X_n-X}>2^{-p})\le2^{-q}.
  \]
\end{enumerate}
Then \(X_n\to X\) in \(L^1\) with a computable modulus, and consequently
\[
 \E[X_n]\longrightarrow\E[X].
\]
\end{theorem}

\begin{proof}
Given output precision \(p\), choose \(K\) so that
\(2\E[G\ind_{G>K}]\le2^{-p-2}\).  Put
\(\eta=2^{-p-2}\).  Choose \(q\) so that
\(2K2^{-q}\le2^{-p-2}\), and then choose
\(n\ge N(p+2,q)\).  The truncation estimate gives
\(\E\abs{X_n-X}\le3\cdot2^{-p-2}<2^{-p}\), after allowing one additional guard
bit for name approximation.  Finally,
\[
 \abs{\E[X_n]-\E[X]}\le\E\abs{X_n-X}\le2^{-p}.
\]
Every choice is primitive recursive from the supplied moduli and rational
bounds.
\end{proof}

\begin{corollary}[Bounded convergence]\label{cor:bounded-convergence}
If \(\abs{X_n},\abs X\le M\) for a fixed rational \(M\) and \(X_n\to X\) in
probability with a modulus, then expectations converge effectively.
\end{corollary}

\begin{proof}
Take \(G=M\).  Its tail above any \(K>M\) is zero, so
\Cref{thm:dct} applies.  Directly,
\(\E\abs{X_n-X}\le\eta+2M\Pp(\abs{X_n-X}>\eta)\).
\end{proof}

\begin{theorem}[Vitali form]\label{thm:vitali}
Effective convergence in probability plus effective uniform integrability of
\((X_n)\) and the limit \(X\) implies effective \(L^1\) convergence and hence
convergence of expectations.
\end{theorem}

\begin{proof}
Use \Cref{lem:truncation} with separate tail bounds for \(X_n\) and \(X\):
\[
 \abs{X_n-X}\le
 \abs{X_n-X}\ind_{\abs{X_n}\le K,\abs X\le K}
 +\abs{X_n}\ind_{\abs{X_n}>K}
 +\abs X\ind_{\abs X>K}.
\]
The first term is bounded by
\(\eta+2K\ind_{\abs{X_n-X}>\eta}\).  Choose \(K,\eta\), and the probability
precision exactly as in \Cref{thm:dct}.
\end{proof}

\subsection{Summary of constructive hypotheses}

The classical convergence theorems remain valid, but their useful
\(\ipra\) forms expose the missing finite data:
\begin{itemize}
\item monotone convergence needs no rate for a lower-real conclusion, but a
  regular-real value needs a tail modulus;
\item Fatou needs effective names for the finite minima and the limiting
  infima;
\item dominated convergence needs an integrable-tail modulus and a convergence
  modulus;
\item exchanging expectation with a signed infinite sum needs absolute
  expectation tails or another uniform-integrability certificate.
\end{itemize}
The moving-spike examples show that these are mathematical requirements, not
merely artifacts of formalization.
